% Reference: https://cp-algorithms.com/graph/edmonds_karp.html
\subsection{网络流基础与最大流最小割}
\label{math:flow-basics}
\subsubsection*{容量、流量与流量守恒}
每条有向边有容量 $c$ 与流量 $f$，满足 $0\le f\le c$；除源点 $s$、汇点 $t$ 外，每点流入等于流出。最大流要求在这些限制下，使源点送到汇点的总流量尽可能大。

\subsubsection*{残量网络与增广}
对原边 $u\to v$，残量图保留容量 $c-f$ 的正向边和容量 $f$ 的反向边。沿正残量的源汇路径增广，增量取路径最小残量；经过反向边相当于撤销之前的部分流量，从而调整已有选择。

本章每次 \InlineCode{add} 都建立独立边对，编号为 $id,id\mathbin{\oplus}1$。反向残量边不是另一条原图边；若输入还含 $v\to u$，应另外加一对。平行边与自环也各自拥有独立的正反边，不通过邻接位置猜反向边。

\subsubsection*{EK 与 Dinic}
EK 每次用 BFS 找边数最少的一条增广路，再沿前驱边回溯推流。Dinic 每轮先用 BFS 分层，再用 DFS 在层次图中寻找阻塞流。两者维护相同的残量边结构，前者适合理解与较小网络，后者通常作为比赛中的通用最大流模板。

\subsubsection*{最大流最小割}
割将顶点分成含源点的 $S$ 与含汇点的 $T$，割容量是所有原始 $S\to T$ 边的容量和。任何可行流都不能超过任意一个割的容量。

当残量网络中再无源汇路径时，从源点沿正残量边可达的点集构成 $S$。原图跨向 $T$ 的边均饱和，反向跨割流量为零，所以割容量等于当前流量；当前流即最大流，该割即最小割。若只是达到人为流量上限而停止，不能据此认定得到最小割。

本章接口的 \InlineCode{flow(s,t,limit)} 均修改残量状态，并返回\textbf{本次额外增广量}；再次调用从当前残量图继续。初始反向容量为零时，原边净流量可由 \InlineCode{e[id^1].cap} 读取。需要重新计算原图最大流时应重新初始化、加边。
